\documentclass[11pt]{article}
\usepackage[a4paper,margin=1in]{geometry}
\usepackage{amsmath,amssymb}
\usepackage{graphicx}
\usepackage{xcolor}
\usepackage{hyperref}
\hypersetup{colorlinks=true,linkcolor=blue,urlcolor=blue,citecolor=blue}
\newcommand{\rt}[1]{\texttt{#1}}
\newcommand{\fn}[1]{\texttt{#1}}
\newcommand{\wcc}{\mathrm{wcc}}
\newcommand{\scc}{\mathrm{scc}}
\newcommand{\ifgraphic}[2]{\IfFileExists{#1}{\includegraphics[width=#2]{#1}}%
  {\fbox{\parbox{#2}{\centering\small figure \texttt{\detokenize{#1}} pending}}}}

\title{\textbf{Report R14 --- one figure for the WCC/SCC planes and the flow
between them}\\
\large Candidate layouts}
\author{Tier 1 analysis}
\date{2026-08-04}

\begin{document}
\maketitle

\noindent R9 spends three figures on a single idea: every rule has a position in
the sector plane, a position in the monitored-attractor plane, and dissipation
moves it from the first to the second. Two candidates follow, then the three
first-pass layouts they replaced.

\medskip
\noindent\textbf{The fact that shapes both.} Of the $230$ rules that move,
$220$ move \emph{straight down} and all $230$ decrease $b$; not one moves
horizontally. That is forced rather than empirical: one attractor per sector
means $n_\scc=n_\wcc$, so $a_\att=a$ and the move can only be vertical. A
horizontal component therefore exists \emph{if and only if} the rule has a
sector holding more than one attractor --- R9~\S7.3's twelve, of which ten show
it (rules $235$ and $249$ have a constant ratio of $2$, which cancels out of a
growth base). So the ten off-axis arrows are not a scatter of outliers; they are
exactly the multi-attractor rules, and both candidates below draw only those.

\vfill
\begin{figure}[htbp]
\centering
\ifgraphic{../../figures/fig_journal_overlay2_obc0.pdf}{0.72\textwidth}
\caption{\textbf{A --- overlay.} One plane. Each rule appears twice, hollow where
its sectors put it and filled where its attractors do. The $220$ forced-vertical
moves are hairline stems with no arrowhead --- they carry no information the
paired markers do not already give --- and only the ten \emph{exceptions} are
drawn as arrows, labelled. Most compact of the layouts, and the only one in which
both coordinates live on the same axes.}
\end{figure}

\clearpage
\begin{figure}[htbp]
\centering
\ifgraphic{../../figures/fig_journal_stack_obc0.pdf}{0.66\textwidth}
\caption{\textbf{B --- stacked, one shared $x$-axis, split $y$.} The
$x$-coordinate means the same thing in both planes, so it is drawn once, at the
bottom. The two $y$-axes are separate, marked by the break slashes on the left
spine, and the gap between the panels is left completely empty --- no spine, no
ticks, no panel titles --- so an arrow never crosses a line or a label. Because
$a_\att=a$ for every rule but ten, almost every arrow is then exactly vertical
and the ten that lean are the multi-attractor rules (orange).}
\end{figure}

\clearpage
\section*{First pass, for the record}

\noindent These three were tried first and are kept only to show what was
rejected: C and D put the two planes in separate coordinate systems, so the eye
cannot compare positions and the verticality of the flow is invisible; E buries
it entirely.

\begin{figure}[htbp]
\centering
\ifgraphic{../../figures/fig_journal_twin_h_obc0.pdf}{0.92\textwidth}
\caption{\textbf{C --- twin panels, horizontal.} Arrows bundle in the gutter and
the gap costs a third of the width.}
\end{figure}

\begin{figure}[htbp]
\centering
\ifgraphic{../../figures/fig_journal_slope_obc0.pdf}{0.36\textwidth}
\caption{\textbf{D --- slopegraph} of the product $ab$ alone. Discards the
$a/b$ split but makes the exclusion curve the picture: everything starts on or
above $2$, the unitary rules run exactly horizontal, and almost everything else
drops below. Still useful as a small companion panel.}
\end{figure}

\begin{figure}[htbp]
\centering
\ifgraphic{../../figures/fig_journal_inset_obc0.pdf}{0.72\textwidth}
\caption{\textbf{E --- inset.} Reflects the asymmetry of the two maps --- the
sector plane is exact for both experiments, the attractor plane is
monitored-only --- but loses the flow altogether.}
\end{figure}

\clearpage
\section*{Recommendation}

\noindent \textbf{A} for a single-panel figure, \textbf{B} when the two planes
must stay visually distinct, and \textbf{D} as a narrow companion to either if
the $ab\ge2$ theorem needs to be shown as well.

A and B encode the same claim in two ways: the flow is one-dimensional except
for ten rules, and those ten are exactly the ones whose sectors carry several
attractors. A says it by superposing the planes; B says it by stacking them on a
common $x$ so that ``vertical'' is literal.

\section*{Reproduce}

\begin{verbatim}
python -m qca_fragmentation.viz.journal --bc obc0
\end{verbatim}

\noindent writes all layouts as
\fn{figures/fig\_journal\_\{style\}\_obc0.\{pdf,png\}}; the module is
\fn{code/qca\_fragmentation/viz/journal.py}, and \fn{journal.off\_axis()}
returns the ten exception rules.

\end{document}
